\selectlanguage{spanish}

\chapter\*{Resumen}

Ubicado en la Pampa Amarilla, en Argentina, el Observatorio Pierre Auger combina diferentes técnicas de detección para estudiar los rayos cósmicos de ultraalta energía mediante un enfoque multihíbrido.
Para extender las mediciones hacia energías más bajas, se instalaron dos arreglos de detectores de superficie más densos, el \sd{750} y el \sd{433}, anidados en arreglo de 1500,m, con separaciones de 750,m y 433,m, respectivamente.
Esta tesis se centra en el \sd{433}, que utiliza detectores Cherenkov de agua para medir el espectro de energía de los rayos cósmicos a partir de $63,\mathrm{PeV}$.
Su rango de energía permite caracterizar el ablandamiento del espectro conocido como la segunda rodilla.
Una caracterización completa de esta estructura aporta información valiosa para comprender la transición de rayos cósmicos galácticos a extragalácticos.

El análisis se basa en \SDFourThreeThreeNEvents{} eventos registrados entre julio de 2018 y diciembre de 2021, con ángulos cenitales menores que $40^\circ$.
Por encima del umbral de energía, el detector puede considerarse eficiente, lo que permite calcular la exposición geométricamente sin recurrir a simulaciones.

Dado el número extremadamente limitado de eventos medidos simultáneamente por el \sd{433} y el detector de fluorescencia, no es posible realizar una calibración directa de la energía.
Sin embargo, los eventos medidos simultáneamente por el \sd{433} y el \sd{750} permiten calibrar el tamaño de la lluvia medido por el \sd{433} usando como referencia la energía del \sd{750}.
De este modo, el \sd{433} hereda la escala de energía del detector de fluorescencia.


Una vez reconstruida la energía del rayo cósmico primario, se obtiene el espectro curdo.
Para medir la segunda rodilla, se realiza un ajuste de verosimilitud en el que el modelo del espectro se convoluciona con la respuesta del detector mediante una matriz de migración para remover los efectos producidos por la eficiencia, la resolucion y el sesgo del detector.

El espectro se describe mediante una ley de potencias con un quiebre suave correspondiente a la segunda rodilla.
Utilizando únicamente datos del \sd{433}, la segunda rodilla se mide en $E\_{01}=(2.8\pm0.5\_{\mathrm{stat}}\pm0.7\_{\mathrm{syst}})\times10^{17},\mathrm{eV}$ y el índice espectral cambia de $\gamma_0=3.04\pm0.02\_{\mathrm{stat}}\pm0.09\_{\mathrm{syst}}$ por debajo del quiebre a $\gamma_1=3.34\pm0.08\_{\mathrm{stat}}\pm0.11\_{\mathrm{syst}}$ por encima.
La incerteza sistemática de la escala de energía del detector de fluorescencia domina la incertidumbre sistemática en el flujo y en la posición de esta estructura, mientras que la incertidumbre de la calibración tiene el mayor efecto sobre los índices espectrales.

Además, se presenta una actualización del espectro del \sd{750} publicado en 2021, utilizando la escala de energía actualizaada del Observatorio.
Usando los eventos registrados por ambos arreglos, se vuelve a estimar la eficiencia de detección del \sd{750}.
Dado que la posición de la segunda rodilla no puede medirse utilizando únicamente datos del \sd{750}, se obtiene un espectro combinado preliminar con datos de ambos arreglos.
Con esta combinación, el objetivo es acotar mejor la posición de la segunda rodilla y extender la medición del espectro más allá del tobillo, contribuyendo así a comprender la transición de rayos cósmicos galácticos a extragalácticos.
